% Part: computability
% Chapter: computability-theory
% Section: total

\documentclass[../../../include/open-logic-section]{subfiles}

\begin{document}

\olfileid{cmp}{thy}{tot}
\olsection{د هرځاے تعريف کېدو مسئله ناپرېکړه کېدونکې ده}

راځئ د~$s$-$m$-$n$ قضيې د کارولو يوه بله بېلګه وګورو، څو وښيو چې
يو څيز محاسبه کېدونکے نۀ دے۔ $\fn{Tot}$ د هرځاے تعريف شويو محاسبه
کېدونکو تابعو د شاخصونو سټ وي، يعنې
\[
\fn{Tot} = \Setabs{x}{\text{د هر $y$ دپاره، $\cfind{x}(y)\fdefined$}}.
\]

\begin{prop}
\ollabel{prop:total}
$\fn{Tot}$ محاسبه کېدونکے نۀ دے۔
\end{prop}

\begin{proof}
د دې د ليدلو دپاره چې~$\fn{Tot}$ محاسبه کېدونکے نۀ دے، دا ښودل بس
دي چې~$K$ ورته راکمېدونکے دے۔ $h(x,y)$ داسې تعريف کړئ:
\[
h(x,y) \simeq
\begin{cases}
0 & \text{که $x \in K$ وي} \\
\fundefined & \text{که داسې نۀ وي۔}
\end{cases}
\]
پام وکړئ چې~$h(x,y)$ بيخي پر~$y$ تکيه نۀ لري۔ دا ليدل بايد ستونزمن
نۀ وي چې~$h$ جزوي محاسبه کېدونکې ده: پر ننوتونو~$x, y$ دا~$h$ په
دې ډول محاسبه کوو چې لومړے تابع~$\cfind{x}$ پر ننوت~$x$ تقليدوو؛
که دا محاسبه ودرېږي، $h(x,y)$، $0$ راګرځوي او درېږي۔ نو $h(x,y)$ يوازې
$\Zero(\umin{s}{T(x,x,s)})$ ده، چېرته چې~$\Zero$ ثابت صفر تابع ده۔
\textbf{د ژباړې د نسخې سمون (OLCMP-035):} سرچينې د دې محاسبې د
بيان په پيل کښې يو ناسمه زياته کلمه لرله؛ طبيعي پښتو جمله د هغې له
نقلولو پرته ژباړل شوې ده۔

د~$s$-$m$-$n$ قضيې په کارولو داسې بنسټيزه بازګشتي تابع~$k(x)$ شته
چې د هر~$x$ او~$y$ دپاره
\[
\cfind{k(x)}(y) =
\begin{cases}
0 & \text{که $x \in K$ وي} \\
\fundefined & \text{که داسې نۀ وي۔}
\end{cases}
\]
نو~$\cfind{k(x)}$ هله هرځاے تعريف شوې ده چې~$x \in K$ وي، او که نه
نو تعريف شوې نۀ ده۔ په دې توګه~$k$ له~$K$ څخه~$\fn{Tot}$ ته راکمونه
ده۔
\end{proof}

\begin{digress}
په حقيقت کښې~$\fn{Tot}$ آن په محاسبوي ډول د شمېر وړ هم نۀ دے؛
پېچلتيا ئې د ``حسابي پوړيز نظام'' په لا لوړه برخه کښې پرته ده۔ خو
دلته به د دې پياوړتيا په اړه اندېښنه نۀ کوو۔
\end{digress}

\end{document}
